\chapter{Conclusion}
\label{chap:conclusion}
The purpose of this project was to model, simulate and control an inverted pendulum mounted on a conveyor-driven cart system. The project combined theoretical modelling, controller design, simulation, PLC implementation and experimental validation of the physical setup. Due to the unstable and nonlinear nature of the inverted pendulum system, the project provided both theoretical and practical challenges related to modern control engineering and automation systems.
Mathematical models of the system were developed using both the Newton-Euler and Euler-Lagrange methods. Although the two approaches were derived from different physical principles, both methods resulted in equivalent dynamic models after linearization around the upright equilibrium position. The linearized models were implemented in MATLAB/Python and used for simulation, controller development and performance analysis.
Several controllers were designed and implemented, including LQR, LQI and cascaded PID. In addition, an energy-based swing up controller was implemented to move the pendulum from the downward hanging position into a range where the balancing controllers could stabilize the system. All controllers were successfully implemented on the physical PLC-based setup and were capable of balancing the pendulum under normal operating conditions.
The experimental results showed that the physical system generally behaved similarly to the simulated models, confirming that the developed mathematical models captured the dominant dynamics of the system. However, noticeable differences between simulation and physical behavior were observed during testing. These differences were primarily caused by unmodelled physical effects such as static friction, backlash, sensor noise, filtering delays actuator saturation and small mechanical imperfections in the setup. In particular, static friction in the conveyor system appeared to have a larger impact on the controller behavior than initially expected, especially for the controllers containing integral action.
Among the tested controllers, the LQR controller provided smooth and stable balancing with relatively low oscillation while the LQI controller improved steady-state position regulation through the addition of integral action. The cascaded PID controller was also able to stabilize the pendulum, although it required substantial manual tuning compared to the state-space based controllers. The project furthermore demonstrated several practical implementation challenges related to real-time control systems including PLC cycle time limitations, encoder quantization, OPC UA communication delays and noisy velocity estimation.
Overall, the project succesfully demonstrated both the theoretical and practical aspects of modelling and controlling an unstable dynamic system. The developed controllers achieved stable balancing of the inverted pendulum and the project provided valuable insight into dynamic modelling, simulation, controller implementation and real-world automation system integration.